\scanpage{197}
Nous allons donner quelques \unsure{rappel\deleted{a}s|rappels} d'Algebre Commutative:

(i) Soit $A$ un anneau \uline{local} \uline{noetherien}. Alors $\dim\,(A)={}$ nombre \unsure{minimmum|minimum}
de generateurs d'un idea\deleted{\unsure{m}}l de definition de $A$. En particulier $\dim(A)$
est finie. Si $A$ \inserted{donc} est un anneau noetherien \unsure{quelconaue|quelconque}, $\underline p$ un ideal premier
de $A$, comme le rang de $\underline p$ est egal evidemment a $\dim A_{\underline p}$, il
s'ensuit qu'il est fini.

(ii) Si $A$ est un anneau \deleted{\unsure{regulier}}\inserted{local} regulier, il satisfait la condition
des chaines (on peut referer a l'article de Serre au Japon). C\uline{'}est donc
vrai aussi par suite pour un anneau regulier quelconque (pas nec.
local), donc pour un quotient d'un tel anneau, i.e. pratiquement pour
\unsure{\deleted{b}tous|bous} les anneaux noetheriens rencontres dans la nature.

(iii) Soient $k$ un corps, $A$ une algebre integre de type fini sur $k$,
$m$ un ideal maximal de $A$, alors $\dim A_m={}$ degre de \unsure{transcendance|transcsndance} de $K$ sur
$k$, ou $K$ est le corps des fractions de $A$. Par suite $A$ est \unsure{equidimansionnel|equidimensionnel}
(et meme strictement, en vertu de (ii)) de \unsure{dimsnsion|dimension} egale au dit
degre de transcendance.

(iv) Enfin, nous utiliserons au Chap.III la definition de la dimension
d'un anneau local noeth. $A$ par le \inserted{degre du} polynome de Hilbert-Samuel,
qu'il\deleted{x} faut donc rappeler ici.

En langage geometrique, (i) signifie que si on a un preschema
\deleted{\unsure{irreductible $X$ et|irreductible $X$}} noetherien $X$ et une partie fermee irreductible $Y$,
alors si $X$ est affine, $\mathrm{codim}_X Y$ est egal au nombre minimum de sections
de $\underline O_X$ telles que $Y$ soit une composante de l'ensemble des zeros communs
a ces sections; et en tous cas ($X$ affine ou non) cette codimension est
finie. (ii) signifie que pratiquement, les preschemas noetheriens
qu'on rencontre satisfont la condition des chaines. Quant a (iii), il
implique qu'un \unsure{presbhama|preschema} algebrique irreductible est strictement
equidimensionnel, et que sa dimension est identique a celle de
toute partie ouverte non vide de $X$ (ce qui ne serait pas vrai par
exemple de $\mathrm{Spec}(V)$, ou $V$ est un anneau de valuation discrete, bien
que $\mathrm{Spec}(V)$ soit manifestement strictement equidimensionnel). \unsure{Enonçnons|Enonçons}

\uline{Proposition} 7. Soit $X$ un preschema algebrique sur le corps $k$, supposons
$X$ irreductible de point generique $z$. Alors $X$ est strictement equidimensionnel
et de dimension egal a $\mathrm{deg.tr.}_k k(z)$ .

\uline{Corollaire} 1 \quad La dim. de $X$ est egale a celle de toute partie ouverte
non vide. \qquad - \quad On en deduit:

\uline{Corollaire} 2 \quad Soit $X$ un \unsure{preschama|preschema} algebrique (pas nec. irred.). Pour
tout point \deleted{ferme} $x$ de $X$, on a $\dim_x X=\mathrm{Sup}\,\dim\,(X_1)$ ou $X_1$ parcourt les
composantes de $X$ contenant $x$. Si $x$ est ferme, on a $\dim_x X=\mathrm{codim}_X x$
% archive sheet number: 196
